\documentclass[11pt,a4paper]{article}
\usepackage[utf8]{inputenc}
\usepackage{graphicx}
\usepackage{amsmath,amssymb}
\usepackage{booktabs}
\usepackage{geometry}
\usepackage{hyperref}
\usepackage{cite}

\geometry{margin=1in}

\title{\textbf{FMHMO-Overlap: Overlapping Hypergraph Community Detection from Fuzzy Memberships}}
\author{\textbf{Piyush Agarwal} \and \textbf{Suraj Kumar Tiwari}}
\date{October 2026}

\begin{document}

\maketitle

\begin{abstract}
Hypergraphs model complex higher-order group relationships where each hyperedge can connect an arbitrary number of nodes. In a recent 2025 study published in \textit{IEEE Transactions on Fuzzy Systems}, Xiao et al. introduced \textbf{FMHMO}, which determines node-to-community fuzzy memberships $u(v, C)$ to assist partition recovery. However, despite the core premise of computing ``fuzzy memberships'', the authors' final step applies a rigid $\operatorname{argmax}$ operation, discarding all multi-community distribution information and producing exclusively crisp partitions. Furthermore, the memberships were never validated on overlapping ground truth. In this work, we develop \textbf{FMHMO-Overlap}, an end-to-end framework that exposes the continuous $n \times c$ membership matrix $U$ and derives overlapping communities via four principled assignment rules (R0 argmax, R1 fixed threshold $\tau$, R2 relative threshold $\rho$, and R3 top-$k$). We design a planted-overlap synthetic hypergraph benchmark and evaluate using overlap-aware metrics: Overlapping NMI, Omega Index, and Overlapping F1. Empirical reproduction on the original paper's datasets replicates reported crisp performance exactly on \texttt{workspace} (NMI 0.836, ARI 0.849, F1 0.924). On planted overlap benchmarks ($p_{\text{ov}} = 0.20$), our framework recovers multi-community nodes with an Overlap-Node F1 of \textbf{0.765} and Omega Index of \textbf{0.828}, whereas the crisp baseline achieves \textbf{0.000}. On crisp controls ($p_{\text{ov}} = 0.0$), our thresholding rules maintain a False Overlap Rate of \textbf{0.00\%}. Finally, we analyze real-world hypergraphs (\texttt{hospital} and \texttt{workspace}) and uncover interpretable cross-department bridge nodes.
\end{abstract}

\section{Introduction}
Hypergraph community detection is pivotal in modeling systems where relationships are inherently group-based, such as scientific co-authorships, cellular protein complexes, and contact networks. Standard graph projections (like clique expansions) often induce severe structural distortion and information loss.

To address this, Xiao et al. (IEEE TFS 2025) proposed the \textbf{FMHMO} framework, which operates in three stages:
\begin{enumerate}
    \item \textbf{HS-HR (Hyperedge Similarity-based Reduction):} Projects hyperedges onto a weighted graph where edge weights $w_{ij}$ combine hyperedge topology and average node similarity parameterized by $\gamma$.
    \item \textbf{Modularity Optimization:} Partitions the reduced graph using Louvain/multilevel modularity maximization.
    \item \textbf{FM-HPR (Partition Recovery):} Maps the partition back to individual nodes. Nodes in multi-community hyperedges (bridge nodes) are resolved by calculating node-to-community fuzzy memberships $u(v, C)$.
\end{enumerate}

\subsection{The Core Research Gap}
Although the paper computes continuous membership grades $u(v, C) \in [0, 1]$, its final recovery step applies:
\begin{equation}
    X[v] = \arg\max_{C} u(v, C)
\end{equation}
All non-maximal membership values are permanently discarded. Consequently, the algorithm produces a purely crisp partition without overlap. Furthermore, the paper never validates memberships on overlapping ground truth and relies exclusively on crisp metrics (standard NMI, ARI) incapable of evaluating multi-membership accuracy.

\section{Proposed Method: FMHMO-Overlap}

\subsection{Continuous Membership Matrix Extraction}
For a hypergraph $G = (V, E)$ with $n = |V|$ nodes and $c$ communities, we extract matrix $U \in \mathbb{R}^{n \times c}$ such that:
\begin{itemize}
    \item For non-bridge nodes whose hyperedges all belong to community $C^*$, the membership row is one-hot: $U[v, C^*] = 1.0$ and $U[v, C] = 0.0$ for $C \neq C^*$.
    \item For bridge nodes, $U[v, :]$ stores the normalized membership grades across neighbor communities $C$:
    \begin{equation}
        u(v, C) = \frac{\operatorname{attr}(v, C)}{\sum_{C' \in \Gamma_C(v)} \operatorname{attr}(v, C')}
    \end{equation}
    where $\operatorname{attr}(v, C) = \sum_{e \in \Gamma_e(v)} u(v, e) \cdot \frac{|e \cap C|}{|e|}$.
\end{itemize}
Every row strictly satisfies $\sum_{j=0}^{c-1} U[v, j] = 1.0$ and $U[v, j] \ge 0$.

\subsection{Assignment Rules}
We implement four modular assignment rules:
\begin{enumerate}
    \item \textbf{R0 (Argmax Baseline):} $v \in C_j \iff j = \arg\max_k U[v, k]$ (Original paper).
    \item \textbf{R1 (Fixed Threshold $\tau$):} $v \in C_j \iff U[v, j] \ge \tau \cup \{\arg\max_k U[v, k]\}$ (Default $\tau = 0.25$).
    \item \textbf{R2 (Relative Threshold $\rho$):} $v \in C_j \iff U[v, j] \ge \rho \cdot \max_k U[v, k]$ (Default $\rho = 0.70$).
    \item \textbf{R3 (Top-$k$):} Assigns node $v$ to its top $k$ communities with $U[v, j] > 0$.
\end{enumerate}

\section{Evaluation Framework}
\subsection{Planted-Overlap Hypergraph Generator}
We generate synthetic hypergraphs with $n$ nodes and $c$ communities. A controlled fraction $p_{\text{ov}} \in [0.0, 0.4]$ of nodes is planted into 2 or 3 communities. Hyperedges of size $d \in [2, 6]$ are drawn from within a community with probability $1 - \xi$, and uniformly at random with noise probability $\xi$. When $p_{\text{ov}} = 0.0$, the benchmark acts as a strictly crisp control.

\subsection{Overlap-Aware Metrics}
We utilize standard overlapping metrics:
\begin{itemize}
    \item \textbf{Overlapping NMI (ONMI):} The LFK / McDaid conditional entropy formulation.
    \item \textbf{Omega Index ($\Omega$):} The overlapping equivalent of the Adjusted Rand Index measuring pairwise co-membership agreement against chance.
    \item \textbf{Overlap-Node Precision, Recall, and F1:} Binary classification performance on identifying which nodes belong to multiple communities.
\end{itemize}

\section{Experimental Results}

\subsection{Baseline Reproduction (E0)}
We reproduced FMHMO on the authors' benchmark datasets (Table \ref{tab:reproduction}). On \texttt{workspace}, our reproduced scores match the published numbers to three decimal places.

\begin{table}[h!]
\centering
\caption{Reproduction comparison with published values (Xiao et al., IEEE TFS 2025).}
\label{tab:reproduction}
\begin{tabular}{lccccc}
\toprule
\textbf{Dataset} & \textbf{Metric} & \textbf{Reproduced (Ours)} & \textbf{Paper Value} & \textbf{Difference} & \textbf{Status} \\
\midrule
\texttt{workspace} & NMI & \textbf{0.836} $\pm$ 0.000 & \textbf{0.836} & 0.000 & Exact Match \\
\texttt{workspace} & ARI & \textbf{0.849} $\pm$ 0.000 & \textbf{0.849} & 0.000 & Exact Match \\
\texttt{workspace} & F1  & \textbf{0.924} $\pm$ 0.000 & \textbf{0.924} & 0.000 & Exact Match \\
\texttt{workspace} & FMI & \textbf{0.891} $\pm$ 0.000 & \textbf{0.891} & 0.000 & Exact Match \\
\texttt{high\_school} & NMI & 0.937 $\pm$ 0.005 & 0.959 & -0.022 & Consistent \\
\texttt{hospital}    & NMI & 0.119 $\pm$ 0.022 & 0.145 & -0.026 & Consistent \\
\bottomrule
\end{tabular}
\end{table}

\subsection{Rule Comparison on Planted Overlap (E1)}
Table \ref{tab:e1} demonstrates the performance of rules as planted overlap increases ($n=200, c=4, \xi=0.20$). While crisp R0 fails completely to detect overlapping nodes (F1 = 0.000), Rule R1 ($\tau=0.25$) achieves \textbf{0.765} Overlap F1 and \textbf{0.828} Omega Index.

\begin{table}[h!]
\centering
\caption{Comparison of assignment rules across planted overlap fractions.}
\label{tab:e1}
\begin{tabular}{llcccc}
\toprule
\textbf{Planted $p_{\text{ov}}$} & \textbf{Rule} & \textbf{ONMI} & \textbf{Omega Index} & \textbf{Overlapping F1} & \textbf{Overlap-Node F1} \\
\midrule
$p_{\text{ov}} = 0.10$ & R0 (Argmax) & 0.779 & 0.816 & 0.939 & 0.000 \\
& \textbf{R1 ($\tau=0.25$)} & \textbf{0.875} & \textbf{0.912} & \textbf{0.973} & \textbf{0.824} \\
& R2 ($\rho=0.70$) & 0.829 & 0.871 & 0.959 & 0.571 \\
\midrule
$p_{\text{ov}} = 0.20$ & R0 (Argmax) & 0.657 & 0.664 & 0.888 & 0.000 \\
& \textbf{R1 ($\tau=0.25$)} & \textbf{0.786} & \textbf{0.828} & \textbf{0.950} & \textbf{0.765} \\
& R2 ($\rho=0.70$) & 0.721 & 0.735 & 0.917 & 0.400 \\
\bottomrule
\end{tabular}
\end{table}

\begin{figure}[h!]
\centering
\includegraphics[width=0.85\textwidth]{E1_rules_vs_pov.png}
\caption{Overlapping NMI and Omega Index vs Planted Overlap fraction ($p_{\text{ov}}$).}
\label{fig:e1}
\end{figure}

\begin{figure}[h!]
\centering
\includegraphics[width=0.95\textwidth]{E2_parameter_sensitivity.png}
\caption{Parameter sensitivity curves for threshold $\tau$ (R1), ratio $\rho$ (R2), and top-$k$ (R3).}
\label{fig:e2}
\end{figure}

\subsection{Control Experiment: Zero False Overlap (E4)}
On purely disjoint hypergraphs ($p_{\text{ov}} = 0.0$), Rules R1 ($\tau \ge 0.25$) and R2 ($\rho \ge 0.70$) yield a False Overlap Rate of \textbf{0.00\%}. This proves that our thresholding mechanism does not hallucinate artificial overlaps when the underlying network is crisp.

\subsection{Real-World Case Studies (E5)}
On real networks, FMHMO-Overlap detects cross-community bridges:
\begin{itemize}
    \item \textbf{\texttt{hospital} ($n=75, m=1825$):} Under R1 ($\tau=0.25$), 36.0\% of nodes exhibit overlap. Node 74 (degree 11) has membership $u(74, C_5) = 0.542$ and $u(74, C_2) = 0.410$, acting as a critical bridge between two hospital staff clusters.
    \item \textbf{\texttt{workspace} ($n=92, m=788$):} Under R1, 10.87\% of nodes overlap. Node 24 (degree 9) connects Community 2 ($u=0.541$) and Community 1 ($u=0.395$).
\end{itemize}

\section{Academic Limitations}
We explicitly document two main limitations:
\begin{enumerate}
    \item \textbf{Single-hyperedge nodes cannot overlap:} Nodes participating in only one hyperedge have one-hot memberships ($u=1.0$) and cannot exhibit multi-membership.
    \item \textbf{Post-partition dependency:} Memberships are calculated after the initial crisp hyperedge partition. If all hyperedges of an overlapping node are clustered into the same community, the overlap cannot be detected.
\end{enumerate}

\section{Conclusion}
FMHMO-Overlap resolves the critical information-loss limitation of the 2025 IEEE TFS paper by exposing the complete fuzzy membership matrix $U$ and applying principled thresholding rules. Our evaluation framework demonstrates significant gains in overlap detection (+12.9\% ONMI at $p_{\text{ov}}=0.2$) while maintaining 0.00\% false overlap on crisp controls.

\begin{thebibliography}{99}
\bibitem{xiao2025} J. Xiao, L. Ma, J. Cao, and X. Xu, ``Hypergraph Community Detection With Fuzzy Memberships,'' \textit{IEEE Trans. Fuzzy Syst.}, vol. 33, no. 9, pp. 3255--3268, 2025.
\bibitem{lancichinetti2009} A. Lancichinetti, S. Fortunato, and J. Kertész, ``Detecting the overlapping and hierarchical community structure in complex networks,'' \textit{New J. Phys.}, vol. 11, no. 3, p. 033015, 2009.
\bibitem{mcdaid2011} A. F. McDaid, D. Greene, and N. Hurley, ``Normalized Mutual Information to evaluate overlapping community findings,'' \textit{arXiv:1110.2515}, 2011.
\bibitem{ruggeri2023} N. Ruggeri, M. Contisciani, F. Battiston, and C. De Bacco, ``Community detection in large hypergraphs,'' \textit{Sci. Adv.}, vol. 9, no. 28, p. eadg9159, 2023.
\end{thebibliography}

\end{document}
